% ============================================================
% TRABAJOS PRÁCTICOS DE TALLER
% ============================================================
\clearpage
\section*{TRABAJOS PRÁCTICOS DE TALLER}
\subsection*{Prácticas formativas de Instrumentos Integrados}

\begin{recuadro}
\textbf{Normas generales:} trabajar con las herramientas indicadas y devolverlas a su lugar;
registrar cada práctica en la planilla; respetar las normas de seguridad y de ESD; entregar las
piezas realizadas (terminales, conectores, mazos) con la tarjeta de identificación del grupo para
su evaluación.
\end{recuadro}

\subsubsection*{TP 1 -- Pelado y crimpado de terminales}
\textbf{Objetivo:} lograr crimpados aprobables y validarlos.\\
\textbf{Materiales:} cable AS22759 de 22, 18 y 14 AWG, peladora de precisión, terminales de ojal y
empalmes aislados de los tres colores, pinza de trinquete, dinamómetro o banco de tracción, lupa.\\
\textbf{Procedimiento:} pelar y crimpar 5 terminales y 2 empalmes por sección; inspeccionar con lupa
según la tabla de la Unidad 7; ensayar a tracción 2 muestras por sección y registrar la fuerza;
analizar los defectos encontrados.

\subsubsection*{TP 2 -- Contactos de conector y D-sub de alta densidad}
\textbf{Objetivo:} armar un conector D-sub con contactos crimpados según un pinout.\\
\textbf{Materiales:} conector DA-26 de alta densidad con contactos de crimpar, crimpadora y
posicionador adecuados, herramientas de inserción y extracción, backshell, cable multiconductor
blindado.\\
\textbf{Procedimiento:} preparar el cable, crimpar los contactos, insertarlos según el pinout,
colocar el backshell y conectar el blindaje; verificar pin a pin con el multímetro; extraer y volver
a insertar un contacto.

\subsubsection*{TP 3 -- Conector circular MIL-DTL-38999}
\textbf{Objetivo:} manejar contactos removibles de conectores militares.\\
\textbf{Procedimiento:} interpretar la designación del conector y su arreglo de inserto; crimpar
contactos tamaño 22D o 20 con la M22520/2-01 o /1-01 y su posicionador; insertarlos con la
herramienta M81969/14 correspondiente; colocar tapones de sellado en las cavidades libres; armar el
backshell; verificar.

\subsubsection*{TP 4 -- Blindajes y manguitos de soldadura}
\textbf{Objetivo:} terminar correctamente el blindaje de un cable de datos.\\
\textbf{Procedimiento:} pelar la cubierta y peinar la malla según la medida; colocar un manguito de
soldadura con cable de drenaje; contraer con pistola de calor y reflector; inspeccionar la fusión
del anillo de estaño; medir continuidad de la malla al drenaje y aislamiento contra los
conductores.

\subsubsection*{TP 5 -- Mazo de cables en tablero}
\textbf{Objetivo:} construir un mazo real a partir de un plano.\\
\textbf{Procedimiento:} dibujar el tablero a escala 1:1; tender los cables identificados; atar con
hilo encerado; armar conectores y terminales; verificar pin a pin con la tabla de cableado; colocar
manguitos y protecciones; registrar el ensayo.

\subsubsection*{TP 6 -- Configuración y calibración}
\textbf{Objetivo:} recorrer el modo de configuración de un equipo integrado.\\
\textbf{Procedimiento:} con un equipo real en banco (o con el simulador o entrenador del
fabricante), ingresar al modo de configuración, cargar las velocidades y arcos de una aeronave dada,
configurar los sensores de motor y simular sus valores; realizar la calibración de nivelación del
AHRS; guardar la configuración y documentar los pasos.

\subsubsection*{TP 7 -- Ensayo pitot-estático}
\textbf{Objetivo:} realizar un ensayo de pérdidas y de indicación.\\
\textbf{Procedimiento:} conectar el banco al sistema de la aeronave o a un tablero didáctico;
realizar el ensayo de pérdidas del sistema estático y del pitot; verificar un altímetro y un
velocímetro en varios puntos; comparar con las tolerancias; completar la tarjeta de trabajo.

\subsubsection*{TP 8 -- Diagnóstico de fallas}
\textbf{Objetivo:} aplicar el método de diagnóstico.\\
\textbf{Procedimiento:} el docente introduce fallas en el tablero didáctico (contacto retraído,
masa abierta, cables A/B invertidos, blindaje cortado, pérdida en una manguera); cada grupo las
diagnostica con el plano, el multímetro y el banco, las corrige y documenta.

\subsubsection*{Formato del informe}
Carátula; objetivo; herramientas y materiales (con números de parte); procedimiento; mediciones y
resultados; fotografías de las piezas realizadas; defectos encontrados y su corrección;
conclusiones sobre la aplicación en el trabajo real.
